%% =====================================================================
%%  5 / 7   Discussion  讨论                                    成员 5
%%  主要任务：解释结果、对照、意义和局限
%%    a. 机制：为什么会出现第 4 部分的结果（不只是复述数字）
%%    b. 对照：与第 2 部分梳理的已有工作一致/不一致在哪，为什么
%%    c. 意义：理论意义与实践意义各一句
%%    d. 局限：数据、规模、评价方式、假设；说清这些局限对结论的威胁
%%  篇幅建议：1.5--2 页；不要引入第 4 部分没有的新结果
%% =====================================================================
\section{Discussion}
\label{sec:discussion}

\todo{PART 5 --- Discussion (Member 5): explain the results, compare them with prior work, and state implications and limitations.}

% ---- 写作骨架 ----
%
% \paragraph{Why the results look like this.}
% 把第 4 部分的现象归因到 2--3 个机制上，每个机制给一句证据（可回引图表）。
%
% \paragraph{Relation to prior work.}
% 与 \ref{sec:litreview} 中的主线逐条对照：哪里一致、哪里相反、差异来源。
%
% \paragraph{Implications.}
% 对后续研究/工程实践各一句话，不要空泛。
%
% \paragraph{Limitations.}
% 明确列出 2--4 条，并说明它们如何限制结论的适用范围。
